\chapter{Anforderungsanalyse}
\label{ch:anforderungsanalyse}

\section{Ist-Analyse}
[PLATZHALTER: Bestehende Situation vor dem Projekt: vorhandene Systeme,
Prozesse, Schwachstellen. Grundlage fuer die Zielsetzung aus
Abschnitt~\ref{sec:projektziel}.]

\section{Lastenheft}
\label{sec:lastenheft}
[PLATZHALTER: Kurzdarstellung der Anforderungen aus Nutzersicht (Was wird
gebraucht und warum?). Vollstaendiges Lastenheft als
Anlage~\ref{anh:lastenheft} beigefuegt -- hier nur Kernanforderungen und
Verweis darauf.]

\section{Soll-Konzept}
[PLATZHALTER: Loesungsidee auf konzeptioneller Ebene, noch ohne technische
Details. Falls mehrere Loesungsalternativen erwogen wurden: hier oder in
Kapitel~\ref{ch:projektplanung} vergleichend darstellen und die Entscheidung
begruenden (Vorgabe: Loesungsalternativen und Entscheidungen sind zu
erlaeutern, nicht nur das Ergebnis).]

\section{Technische Grundlagen}
[PLATZHALTER: Fuer das Verstaendnis noetige Technologien/Konzepte kurz
erlaeutern (nur was fuer Kapitel~\ref{ch:umsetzung} vorausgesetzt wird).
Nicht allgemeingueltige Fachbegriffe/Abkuerzungen hier erklaeren oder ins
Abkuerzungsverzeichnis aufnehmen.]

\section{Pflichtenheft}
\label{sec:pflichtenheft}
[PLATZHALTER: Aus dem Lastenheft abgeleitete technische Umsetzungsvorgaben
(Wie wird es umgesetzt?). Vollstaendiges Pflichtenheft als
Anlage~\ref{anh:pflichtenheft} beigefuegt -- hier nur Kernpunkte und Verweis
darauf.]
